% ---
% file: docs/latex/formulario_basico.tex
% description: Compendio didáctico de fórmulas matemáticas de nivel básico (aritmética, álgebra elemental, trigonometría).
% type: doc/manual
% version: 1.0.0
% date: 2026-08-24
% covers:
%   - src/data/symbols.py
% relations:
%   - docs/mates/formulario_mates.md
%   - docs/fisica/formulario_basico.md
%   - docs/ARCHITECTURE.md
% keywords:
%   - formulas-basicas-mates
%   - aritmetica-basica
%   - algebra-elemental
%   - trigonometria-basica
% ---

\documentclass[11pt,a4paper]{article}
\usepackage[utf8]{inputenc}
\usepackage[T1]{fontenc}
\usepackage[spanish,es-nodecimaldot,es-noquoting]{babel}
\usepackage{amsmath,amssymb,amsfonts}

% Configuración de márgenes y espaciado estándar
\setlength{\textwidth}{16.5cm}
\setlength{\textheight}{24cm}
\setlength{\oddsidemargin}{-0.25cm}
\setlength{\evensidemargin}{-0.25cm}
\setlength{\topmargin}{-1cm}
\setlength{\parindent}{0pt}
\setlength{\parskip}{5pt}

\begin{document}

\section*{Compendio I: Aritmética, Álgebra y Trigonometría}
\addcontentsline{toc}{section}{Compendio I: Aritmética, Álgebra y Trigonometría}

\textit{(Nivel Básico)}

\section*{1. Aritmética}
\addcontentsline{toc}{section}{1. Aritmética}

\subsection*{1.1 Propiedades de la Adición y Multiplicación}
\addcontentsline{toc}{subsection}{1.1 Propiedades de la Adición y Multiplicación}

\begin{itemize}
  \item \textbf{Clausura / Cerradura (Adición):}
  \[
    \forall a, b \in \mathbb{R}: \quad a + b \in \mathbb{R}
  \]
  \item \textbf{Clausura / Cerradura (Multiplicación):}
  \[
    \forall a, b \in \mathbb{R}: \quad a \cdot b \in \mathbb{R}
  \]
  \item \textbf{Conmutativa (Adición):}
  \[
    a + b = b + a
  \]
  \item \textbf{Conmutativa (Multiplicación):}
  \[
    a \cdot b = b \cdot a
  \]
  \item \textbf{Asociativa (Adición):}
  \[
    (a + b) + c = a + (b + c)
  \]
  \item \textbf{Asociativa (Multiplicación):}
  \[
    (a \cdot b) \cdot c = a \cdot (b \cdot c)
  \]
  \item \textbf{Elemento Neutro (Adición):}
  \[
    a + 0 = a
  \]
  \item \textbf{Elemento Neutro (Multiplicación):}
  \[
    a \cdot 1 = a
  \]
  \item \textbf{Elemento Inverso (Adición):}
  \[
    a + (-a) = 0
  \]
  \item \textbf{Elemento Inverso (Multiplicación):}
  \[
    a \cdot \left(\frac{1}{a}\right) = 1 \quad (a \ne 0)
  \]
  \item \textbf{Distributiva:}
  \[
    a \cdot (b + c) = a \cdot b + a \cdot c
  \]
\end{itemize}

\subsection*{1.2 Operaciones con Fracciones}
\addcontentsline{toc}{subsection}{1.2 Operaciones con Fracciones}

\begin{itemize}
  \item \textbf{Suma/Resta con mismo denominador:}
  \[
    \frac{a}{c} \pm \frac{b}{c} = \frac{a \pm b}{c} \quad (c \ne 0)
  \]
  \item \textbf{Suma/Resta con distinto denominador:}
  \[
    \frac{a}{b} \pm \frac{c}{d} = \frac{ad \pm bc}{bd} \quad (b, d \ne 0)
  \]
  \item \textbf{Multiplicación:}
  \[
    \frac{a}{b} \cdot \frac{c}{d} = \frac{a \cdot c}{b \cdot d} \quad (b, d \ne 0)
  \]
  \item \textbf{División:}
  \[
    \frac{\frac{a}{b}}{\frac{c}{d}} = \frac{a \cdot d}{b \cdot c} \quad (b, c, d \ne 0)
  \]
\end{itemize}

\subsection*{1.3 Porcentajes y Proporcionalidad}
\addcontentsline{toc}{subsection}{1.3 Porcentajes y Proporcionalidad}

\begin{itemize}
  \item \textbf{Tanto por ciento ($P = \text{porcentaje}$, $C = \text{cantidad}$, $r = \text{tasa}$):}
  \[
    P = \frac{C \cdot r}{100}
  \]
  \item \textbf{Regla de tres simple directa ($\frac{a}{b} = \frac{c}{x}$):}
  \[
    x = \frac{b \cdot c}{a}
  \]
\end{itemize}

\section*{2. Álgebra}
\addcontentsline{toc}{section}{2. Álgebra}

\subsection*{2.1 Leyes de Exponentes y Radicales}
\addcontentsline{toc}{subsection}{2.1 Leyes de Exponentes y Radicales}

\begin{itemize}
  \item \textbf{Multiplicación de bases iguales:}
  \[
    x^a \cdot x^b = x^{a + b}
  \]
  \item \textbf{División de bases iguales:}
  \[
    \frac{x^a}{x^b} = x^{a - b}
  \]
  \item \textbf{Potencia de una potencia:}
  \[
    (x^a)^b = x^{a \cdot b}
  \]
  \item \textbf{Potencia de un producto:}
  \[
    (x \cdot y)^a = x^a \cdot y^a
  \]
  \item \textbf{Exponente negativo ($x \ne 0$):}
  \[
    x^{-a} = \frac{1}{x^a}
  \]
  \item \textbf{Exponente fraccionario:}
  \[
    x^{\frac{a}{b}} = \sqrt[b]{x^a}
  \]
  \item \textbf{Raíz de un producto:}
  \[
    \sqrt[n]{x \cdot y} = \sqrt[n]{x} \cdot \sqrt[n]{y}
  \]
\end{itemize}

\subsection*{2.2 Productos Notables y Factorización}
\addcontentsline{toc}{subsection}{2.2 Productos Notables y Factorización}

\begin{itemize}
  \item \textbf{Binomio al cuadrado:}
  \[
    (a \pm b)^2 = a^2 \pm 2ab + b^2
  \]
  \item \textbf{Diferencia de cuadrados:}
  \[
    a^2 - b^2 = (a + b)(a - b)
  \]
  \item \textbf{Binomio al cubo:}
  \[
    (a \pm b)^3 = a^3 \pm 3a^2b + 3ab^2 \pm b^3
  \]
  \item \textbf{Suma de cubos:}
  \[
    a^3 + b^3 = (a + b)(a^2 - ab + b^2)
  \]
  \item \textbf{Diferencia de cubos:}
  \[
    a^3 - b^3 = (a - b)(a^2 + ab + b^2)
  \]
  \item \textbf{Trinomio al cuadrado:}
  \[
    (a + b + c)^2 = a^2 + b^2 + c^2 + 2(ab + ac + bc)
  \]
\end{itemize}

\subsection*{2.3 Ecuaciones Lineales}
\addcontentsline{toc}{subsection}{2.3 Ecuaciones Lineales}

\begin{itemize}
  \item \textbf{Forma general ($a \ne 0$):}
  \[
    ax + b = 0 \implies x = -\frac{b}{a}
  \]
\end{itemize}

\section*{3. Trigonometría}
\addcontentsline{toc}{section}{3. Trigonometría}

\subsection*{3.1 Razones Trigonométricas en el Triángulo Rectángulo}
\addcontentsline{toc}{subsection}{3.1 Razones Trigonométricas en el Triángulo Rectángulo}

\begin{itemize}
  \item \textbf{Seno:}
  \[
    \sin(\theta) = \frac{\text{Cateto Opuesto}}{\text{Hipotenusa}}
  \]
  \item \textbf{Coseno:}
  \[
    \cos(\theta) = \frac{\text{Cateto Adyacente}}{\text{Hipotenusa}}
  \]
  \item \textbf{Tangente:}
  \[
    \tan(\theta) = \frac{\text{Cateto Opuesto}}{\text{Cateto Adyacente}}
  \]
  \item \textbf{Cosecante:}
  \[
    \csc(\theta) = \frac{\text{Hipotenusa}}{\text{Cateto Opuesto}} = \frac{1}{\sin(\theta)}
  \]
  \item \textbf{Secante:}
  \[
    \sec(\theta) = \frac{\text{Hipotenusa}}{\text{Cateto Adyacente}} = \frac{1}{\cos(\theta)}
  \]
  \item \textbf{Cotangente:}
  \[
    \cot(\theta) = \frac{\text{Cateto Adyacente}}{\text{Cateto Opuesto}} = \frac{1}{\tan(\theta)}
  \]
\end{itemize}

\subsection*{3.2 Identidades Trigonométricas Fundamentales}
\addcontentsline{toc}{subsection}{3.2 Identidades Trigonométricas Fundamentales}

\begin{itemize}
  \item \textbf{Tangente por cociente:}
  \[
    \tan(\theta) = \frac{\sin(\theta)}{\cos(\theta)}
  \]
  \item \textbf{Cotangente por cociente:}
  \[
    \cot(\theta) = \frac{\cos(\theta)}{\sin(\theta)}
  \]
  \item \textbf{Identidad Pitagórica principal:}
  \[
    \sin^2(\theta) + \cos^2(\theta) = 1
  \]
  \item \textbf{Identidad Pitagórica (Tangente y Secante):}
  \[
    1 + \tan^2(\theta) = \sec^2(\theta)
  \]
  \item \textbf{Identidad Pitagórica (Cotangente y Cosecante):}
  \[
    1 + \cot^2(\theta) = \csc^2(\theta)
  \]
\end{itemize}

\section*{Compendio II: Geometría, Geometría Analítica y Álgebra Lineal}
\addcontentsline{toc}{section}{Compendio II: Geometría, Geometría Analítica y Álgebra Lineal}

\textit{(Nivel Básico)}

\section*{4. Geometría (Plana y del Espacio / Euclidiana)}
\addcontentsline{toc}{section}{4. Geometría (Plana y del Espacio / Euclidiana)}

\subsection*{4.1 Perímetros y Áreas de Figuras Planas (2D)}
\addcontentsline{toc}{subsection}{4.1 Perímetros y Áreas de Figuras Planas (2D)}

\begin{itemize}
  \item \textbf{Perímetro de Triángulo (General):}
  \[
    P = a + b + c
  \]
  \item \textbf{Área de Triángulo (General):}
  \[
    A = \frac{b \cdot h}{2}
  \]
  \item \textbf{Área de Triángulo Equilátero:}
  \[
    A = \frac{\sqrt{3}}{4} l^2
  \]
  \item \textbf{Altura de Triángulo Equilátero:}
  \[
    h = \frac{\sqrt{3}}{2} l
  \]
  \item \textbf{Perímetro de Cuadrado:}
  \[
    P = 4l
  \]
  \item \textbf{Área de Cuadrado:}
  \[
    A = l^2 = \frac{d^2}{2}
  \]
  \item \textbf{Perímetro de Rectángulo:}
  \[
    P = 2(b + h)
  \]
  \item \textbf{Área de Rectángulo:}
  \[
    A = b \cdot h
  \]
  \item \textbf{Paralelogramo:}
  \[
    A = b \cdot h
  \]
  \item \textbf{Perímetro de Rombo:}
  \[
    P = 4l
  \]
  \item \textbf{Área de Rombo:}
  \[
    A = \frac{D \cdot d}{2}
  \]
  \item \textbf{Trapecio:}
  \[
    A = \frac{(B + b) \cdot h}{2}
  \]
  \item \textbf{Perímetro de Polígono Regular ($n$ lados):}
  \[
    P = n \cdot l
  \]
  \item \textbf{Área de Polígono Regular ($n$ lados):}
  \[
    A = \frac{P \cdot a_p}{2}
  \]
  \item \textbf{Circunferencia de Círculo:}
  \[
    C = 2\pi r
  \]
  \item \textbf{Área de Círculo:}
  \[
    A = \pi r^2
  \]
  \item \textbf{Longitud de Arco de Sector Circular:}
  \[
    s = r\theta \quad (\theta \text{ en rad})
  \]
  \item \textbf{Área de Sector Circular:}
  \[
    A = \frac{1}{2} r^2 \theta \quad (\theta \text{ en rad})
  \]
  \item \textbf{Corona Circular:}
  \[
    A = \pi(R^2 - r^2)
  \]
\end{itemize}

\subsection*{4.2 Propiedades Angulares y Polígonos}
\addcontentsline{toc}{subsection}{4.2 Propiedades Angulares y Polígonos}

\begin{itemize}
  \item \textbf{Teorema de Pitágoras (Triángulos rectángulos):}
  \[
    a^2 + b^2 = c^2
  \]
  \item \textbf{Suma de ángulos internos de un triángulo:}
  \[
    \alpha + \beta + \gamma = 180^\circ \quad (\text{o } \pi\text{ rad})
  \]
  \item \textbf{Suma de ángulos internos de un polígono de $n$ lados:}
  \[
    S_{\text{int}} = (n - 2) \cdot 180^\circ
  \]
  \item \textbf{Ángulo interior de un polígono regular:}
  \[
    \theta_{\text{int}} = \frac{(n - 2) \cdot 180^\circ}{n}
  \]
  \item \textbf{Número total de diagonales en un polígono:}
  \[
    N_d = \frac{n(n - 3)}{2}
  \]
\end{itemize}

\section*{5. Geometría Analítica}
\addcontentsline{toc}{section}{5. Geometría Analítica}

\subsection*{5.1 Plano Cartesiano (2D)}
\addcontentsline{toc}{subsection}{5.1 Plano Cartesiano (2D)}

\begin{itemize}
  \item \textbf{Distancia entre dos puntos $P_1(x_1, y_1)$ y $P_2(x_2, y_2)$:}
  \[
    d = \sqrt{(x_2 - x_1)^2 + (y_2 - y_1)^2}
  \]
  \item \textbf{Punto medio $M$:}
  \[
    M = \left(\frac{x_1 + x_2}{2}, \frac{y_1 + y_2}{2}\right)
  \]
  \item \textbf{Coordenada x en división de segmento en una razón $r$ ($\frac{AP}{PB} = r$):}
  \[
    x = \frac{x_1 + r x_2}{1 + r}
  \]
  \item \textbf{Coordenada y en división de segmento en una razón $r$ ($\frac{AP}{PB} = r$):}
  \[
    y = \frac{y_1 + r y_2}{1 + r}
  \]
\end{itemize}

\subsection*{5.2 La Línea Recta en 2D}
\addcontentsline{toc}{subsection}{5.2 La Línea Recta en 2D}

\begin{itemize}
  \item \textbf{Pendiente:}
  \[
    m = \frac{y_2 - y_1}{x_2 - x_1} = \tan(\theta)
  \]
  \item \textbf{Ecuación Punto-Pendiente:}
  \[
    y - y_1 = m(x - x_1)
  \]
  \item \textbf{Ecuación Pendiente-Ordenada al origen:}
  \[
    y = mx + b
  \]
  \item \textbf{Ecuación General ($B \ne 0 \implies m = -\frac{A}{B}$):}
  \[
    Ax + By + C = 0
  \]
  \item \textbf{Ecuación Simétrica / Canónica:}
  \[
    \frac{x}{a} + \frac{y}{b} = 1 \quad (\text{Intersecciones en } (a,0) \text{ y } (0,b))
  \]
  \item \textbf{Rectas Paralelas:}
  \[
    m_1 = m_2
  \]
  \item \textbf{Rectas Perpendiculares:}
  \[
    m_1 \cdot m_2 = -1 \iff m_1 = -\frac{1}{m_2}
  \]
  \item \textbf{Ángulo entre dos rectas:}
  \[
    \tan(\theta) = \left|\frac{m_2 - m_1}{1 + m_1 \cdot m_2}\right|
  \]
\end{itemize}

\section*{6. Álgebra Lineal}
\addcontentsline{toc}{section}{6. Álgebra Lineal}

\subsection*{6.1 Vectores en $\mathbb{R}^n$}
\addcontentsline{toc}{subsection}{6.1 Vectores en $\mathbb{R}^n$}

\begin{itemize}
  \item \textbf{Suma:}
  \[
    \mathbf{u} + \mathbf{v} = (u_1 + v_1, u_2 + v_2, \dots, u_n + v_n)
  \]
  \item \textbf{Producto Escalar por Vector:}
  \[
    c\mathbf{u} = (cu_1, cu_2, \dots, cu_n)
  \]
  \item \textbf{Norma / Módulo Euclidiano ($L_2$):}
  \[
    \|\mathbf{u}\| = \sqrt{u_1^2 + u_2^2 + \dots + u_n^2} = \sqrt{\mathbf{u} \cdot \mathbf{u}}
  \]
  \item \textbf{Vector Unitario:}
  \[
    \hat{\mathbf{u}} = \frac{\mathbf{u}}{\|\mathbf{u}\|}
  \]
\end{itemize}

\subsection*{6.2 Productos Vectoriales}
\addcontentsline{toc}{subsection}{6.2 Productos Vectoriales}

\begin{itemize}
  \item \textbf{Producto Punto (Escalar):}
  \[
    \mathbf{u} \cdot \mathbf{v} = \sum_{i=1}^n u_i v_i = \|\mathbf{u}\| \|\mathbf{v}\| \cos(\theta)
  \]
  \item \textbf{Ortogonalidad:}
  \[
    \mathbf{u} \perp \mathbf{v} \iff \mathbf{u} \cdot \mathbf{v} = 0
  \]
  \item \textbf{Producto Cruz (Vectorial en $\mathbb{R}^3$):}
  \[
    \mathbf{u} \times \mathbf{v} = \det\begin{pmatrix} \mathbf{i} & \mathbf{j} & \mathbf{k} \\ u_1 & u_2 & u_3 \\ v_1 & v_2 & v_3 \end{pmatrix}
  \]
  \item \textbf{Magnitud del Producto Cruz:}
  \[
    \|\mathbf{u} \times \mathbf{v}\| = \|\mathbf{u}\| \|\mathbf{v}\| \sin(\theta)
  \]
  \item \textbf{Producto Triple Escalar (Volumen del paralelepípedo):}
  \[
    [\mathbf{u}, \mathbf{v}, \mathbf{w}] = \mathbf{u} \cdot (\mathbf{v} \times \mathbf{w}) = \det\begin{pmatrix} u_1 & u_2 & u_3 \\ v_1 & v_2 & v_3 \\ w_1 & w_2 & w_3 \end{pmatrix}
  \]
\end{itemize}

\subsection*{6.3 Sistemas de Ecuaciones Lineales}
\addcontentsline{toc}{subsection}{6.3 Sistemas de Ecuaciones Lineales}

\begin{itemize}
  \item \textbf{Sistema Matricial:}
  \[
    A\mathbf{x} = \mathbf{b}
  \]
  \item \textbf{Regla de Cramer ($\det(A) \ne 0$):}
  \[
    x_i = \frac{\det(A_i)}{\det(A)}
  \]
\end{itemize}

\section*{Compendio III: Cálculo Diferencial, Integral y Vectorial}
\addcontentsline{toc}{section}{Compendio III: Cálculo Diferencial, Integral y Vectorial}

\textit{(Nivel Básico)}

\section*{7. Cálculo Diferencial}
\addcontentsline{toc}{section}{7. Cálculo Diferencial}

\subsection*{7.1 Definición de Límite y Continuidad}
\addcontentsline{toc}{subsection}{7.1 Definición de Límite y Continuidad}

\begin{itemize}
  \item \textbf{Definición épsilon-delta:}
  \[
    \lim_{x \to a} f(x) = L \iff \forall \varepsilon > 0, \, \exists \delta > 0 : 0 < |x - a| < \delta \implies |f(x) - L| < \varepsilon
  \]
  \item \textbf{Continuidad en un punto $a$:}
  \[
    \lim_{x \to a} f(x) = f(a)
  \]
\end{itemize}

\subsection*{7.2 Definición Formal de la Derivada}
\addcontentsline{toc}{subsection}{7.2 Definición Formal de la Derivada}

\begin{itemize}
  \item \textbf{Por límite (Cociente diferencial):}
  \[
    f'(x) = \frac{dy}{dx} = \lim_{h \to 0} \frac{f(x + h) - f(x)}{h}
  \]
  \item \textbf{En un punto $a$:}
  \[
    f'(a) = \lim_{x \to a} \frac{f(x) - f(a)}{x - a}
  \]
\end{itemize}

\subsection*{7.3 Reglas y Derivadas Algebraicas Elementales}
\addcontentsline{toc}{subsection}{7.3 Reglas y Derivadas Algebraicas Elementales}

\begin{itemize}
  \item \textbf{Constante:}
  \[
    \frac{d}{dx}(c) = 0
  \]
  \item \textbf{Identidad:}
  \[
    \frac{d}{dx}(x) = 1
  \]
  \item \textbf{Regla de la Potencia:}
  \[
    \frac{d}{dx}(x^n) = n x^{n - 1}
  \]
  \item \textbf{Múltiplo Constante:}
  \[
    \frac{d}{dx}[c \cdot f(x)] = c f'(x)
  \]
  \item \textbf{Suma y Resta:}
  \[
    \frac{d}{dx}[f(x) \pm g(x)] = f'(x) \pm g'(x)
  \]
  \item \textbf{Raíz cuadrada:}
  \[
    \frac{d}{dx}(\sqrt{x}) = \frac{1}{2\sqrt{x}}
  \]
  \item \textbf{Inverso:}
  \[
    \frac{d}{dx}\left(\frac{1}{x}\right) = -\frac{1}{x^2}
  \]
\end{itemize}

\section*{8. Cálculo Integral}
\addcontentsline{toc}{section}{8. Cálculo Integral}

\subsection*{8.1 Definición y Propiedades Fundamentales}
\addcontentsline{toc}{subsection}{8.1 Definición y Propiedades Fundamentales}

\begin{itemize}
  \item \textbf{Antiderivada:}
  \[
    \int f(x) \, dx = F(x) + C \iff F'(x) = f(x)
  \]
  \item \textbf{Linealidad:}
  \[
    \int [a f(x) + b g(x)] \, dx = a \int f(x) \, dx + b \int g(x) \, dx
  \]
\end{itemize}

\subsection*{8.2 Integrales Indefinidas Elementales / Inmediatas}
\addcontentsline{toc}{subsection}{8.2 Integrales Indefinidas Elementales / Inmediatas}

\begin{itemize}
  \item \textbf{Potencias ($n \ne -1$):}
  \[
    \int x^n \, dx = \frac{x^{n+1}}{n+1} + C
  \]
  \item \textbf{Logarítmica:}
  \[
    \int \frac{1}{x} \, dx = \ln|x| + C
  \]
  \item \textbf{Integral de exponencial natural:}
  \[
    \int e^x \, dx = e^x + C
  \]
  \item \textbf{Integral de exponencial general:}
  \[
    \int a^x \, dx = \frac{a^x}{\ln(a)} + C
  \]
  \item \textbf{Integral de Seno:}
  \[
    \int \sin(x) \, dx = -\cos(x) + C
  \]
  \item \textbf{Integral de Coseno:}
  \[
    \int \cos(x) \, dx = \sin(x) + C
  \]
  \item \textbf{Integral de Secante cuadrada:}
  \[
    \int \sec^2(x) \, dx = \tan(x) + C
  \]
  \item \textbf{Integral de Cosecante cuadrada:}
  \[
    \int \csc^2(x) \, dx = -\cot(x) + C
  \]
  \item \textbf{Integral de Secante por Tangente:}
  \[
    \int \sec(x)\tan(x) \, dx = \sec(x) + C
  \]
  \item \textbf{Integral de Cosecante por Cotangente:}
  \[
    \int \csc(x)\cot(x) \, dx = -\csc(x) + C
  \]
\end{itemize}

\section*{9. Cálculo Vectorial y Multivariable}
\addcontentsline{toc}{section}{9. Cálculo Vectorial y Multivariable}

\subsection*{9.1 Curvas Paramétricas y Funciones Vectoriales en $\mathbb{R}^3$}
\addcontentsline{toc}{subsection}{9.1 Curvas Paramétricas y Funciones Vectoriales en $\mathbb{R}^3$}

\begin{itemize}
  \item \textbf{Vector de Posición:}
  \[
    \mathbf{r}(t) = x(t)\mathbf{i} + y(t)\mathbf{j} + z(t)\mathbf{k}
  \]
  \item \textbf{Velocidad:}
  \[
    \mathbf{v}(t) = \mathbf{r}'(t) = x'(t)\mathbf{i} + y'(t)\mathbf{j} + z'(t)\mathbf{k}
  \]
  \item \textbf{Rapidez:}
  \[
    \|\mathbf{v}(t)\| = \frac{ds}{dt} = \sqrt{[x'(t)]^2 + [y'(t)]^2 + [z'(t)]^2}
  \]
  \item \textbf{Aceleración:}
  \[
    \mathbf{a}(t) = \mathbf{r}''(t) = \mathbf{v}'(t)
  \]
\end{itemize}

\subsection*{9.2 Derivadas Parciales de Funciones Escalares $f(x, y, z)$}
\addcontentsline{toc}{subsection}{9.2 Derivadas Parciales de Funciones Escalares $f(x, y, z)$}

\begin{itemize}
  \item \textbf{Derivada parcial respecto a $x$:}
  \[
    \frac{\partial f}{\partial x} = \lim_{h \to 0} \frac{f(x + h, y, z) - f(x, y, z)}{h}
  \]
  \item \textbf{Teorema de Schwarz / Clairaut (Derivadas cruzadas):}
  \[
    \frac{\partial^2 f}{\partial x \partial y} = \frac{\partial^2 f}{\partial y \partial x}
  \]
\end{itemize}

\subsection*{9.3 Gradiente}
\addcontentsline{toc}{subsection}{9.3 Gradiente}

\begin{itemize}
  \item \textbf{Vector Gradiente:}
  \[
    \nabla f = \operatorname{grad}(f) = \frac{\partial f}{\partial x}\mathbf{i} + \frac{\partial f}{\partial y}\mathbf{j} + \frac{\partial f}{\partial z}\mathbf{k}
  \]
\end{itemize}

\section*{Compendio IV: Ecuaciones Diferenciales y Métodos Numéricos}
\addcontentsline{toc}{section}{Compendio IV: Ecuaciones Diferenciales y Métodos Numéricos}

\textit{(Nivel Básico)}

\section*{10. Ecuaciones Diferenciales}
\addcontentsline{toc}{section}{10. Ecuaciones Diferenciales}

\subsection*{10.1 Definiciones y Clasificación}
\addcontentsline{toc}{subsection}{10.1 Definiciones y Clasificación}

\begin{itemize}
  \item \textbf{Problema de Valor Inicial (PVI) (1):}
  \[
    y' = f(x, y)
  \]
  \item \textbf{Problema de Valor Inicial (PVI) (2):}
  \[
    y(x_0) = y_0
  \]
  \item \textbf{Teorema de Existencia:}
  \[
    y' = f(x, y)
  \]
  \item \textbf{Unicidad de Picard-Lindelöf:}
  \[
    y(x_0) = y_0 \implies \exists ! \, y(x) \quad \text{en } I
  \]
\end{itemize}

\subsection*{10.2 Ecuaciones Diferenciales Ordinarias (EDO) de 1er Orden Elementales}
\addcontentsline{toc}{subsection}{10.2 Ecuaciones Diferenciales Ordinarias (EDO) de 1er Orden Elementales}

\begin{itemize}
  \item \textbf{Variables Separables:}
  \[
    g(y) \, dy = f(x) \, dx \implies \int g(y) \, dy = \int f(x) \, dx + C
  \]
  \item \textbf{Forma estándar de Ecuación Lineal de 1er Orden:}
  \[
    y' + P(x)y = Q(x)
  \]
  \item \textbf{Factor Integrante de Leibniz:}
  \[
    \mu(x) = e^{\int P(x) \, dx}
  \]
  \item \textbf{Solución de Ecuación Lineal de 1er Orden:}
  \[
    y(x) = \frac{1}{\mu(x)} \left[ \int \mu(x) Q(x) \, dx + C \right]
  \]
  \item \textbf{Forma de Ecuación Diferencial Exacta:}
  \[
    M(x, y) \, dx + N(x, y) \, dy = 0
  \]
  \item \textbf{Condición y Solución de Ecuación Exacta:}
  \[
    \frac{\partial M}{\partial y} = \frac{\partial N}{\partial x} \implies \Psi(x, y) = C
  \]
  \item \textbf{Factor Integrante (dependiente de x) para Ecuaciones no Exactas:}
  \[
    \mu(x) = e^{\int \frac{M_y - N_x}{N} \, dx}
  \]
  \item \textbf{Factor Integrante (dependiente de y) para Ecuaciones no Exactas:}
  \[
    \mu(y) = e^{\int \frac{N_x - M_y}{M} \, dy}
  \]
\end{itemize}

\section*{11. Métodos Numéricos}
\addcontentsline{toc}{section}{11. Métodos Numéricos}

\subsection*{11.1 Teoría de Errores}
\addcontentsline{toc}{subsection}{11.1 Teoría de Errores}

\begin{itemize}
  \item \textbf{Error Absoluto:}
  \[
    E_a = |x - \hat{x}|
  \]
  \item \textbf{Error Relativo:}
  \[
    E_r = \frac{|x - \hat{x}|}{|x|} \quad (x \ne 0)
  \]
  \item \textbf{Error Relativo Porcentual:}
  \[
    E_p = E_r \cdot 100\%
  \]
  \item \textbf{Error Aproximado:}
  \[
    \varepsilon_a = \left|\frac{x_{\text{actual}} - x_{\text{anterior}}}{x_{\text{actual}}}\right| \cdot 100\%
  \]
\end{itemize}

\subsection*{11.2 Raíces de Ecuaciones No Lineales (Métodos de Intervalo / Cerrados)}
\addcontentsline{toc}{subsection}{11.2 Raíces de Ecuaciones No Lineales (Métodos de Intervalo / Cerrados)}

\begin{itemize}
  \item \textbf{Teorema de Bolzano:}
  \[
    f(a) \cdot f(b) < 0 \implies \exists c \in (a, b) : f(c) = 0
  \]
  \item \textbf{Punto medio en Método de Bisección:}
  \[
    c_n = \frac{a_n + b_n}{2}
  \]
  \item \textbf{Cota de Error en Método de Bisección:}
  \[
    E_n = |r - c_n| \le \frac{b - a}{2^{n+1}}
  \]
  \item \textbf{Número de iteraciones en Método de Bisección:}
  \[
    n \ge \left\lceil \log_2\left(\frac{b - a}{2\varepsilon}\right) \right\rceil
  \]
  \item \textbf{Método de Falsa Posición (Regula Falsi):}
  \[
    c = b - \frac{f(b)(a - b)}{f(a) - f(b)} = \frac{a f(b) - b f(a)}{f(b) - f(a)}
  \]
\end{itemize}


\end{document}
